\section{偏见：比你想象的更微妙}

\subsection{SQL 查询中的性别偏见}

讲者分享了一个亲身经历的案例来说明 AI 偏见的微妙性。他在做一个家谱学项目时，向某个聊天机器人请求生成一段 SQL 查询来查找所有祖先。返回的代码中有一个关键错误：在递归查询中，某一行\textbf{只检查了 father 是否为祖先，而忽略了 mother}。虽然代码的其他部分处理了母亲一侧，但这个特定的递归步骤遗漏了。

\begin{warningbox}{Bug 还是 Bias？一个难以回答的问题}
讲者与妻子讨论了这个问题：这到底是一个\textbf{编程错误}（Bug）还是\textbf{性别偏见}（Gender Bias）？如果训练数据中的家谱代码示例大多以父系为主，模型就会``学到''这种偏好模式。关键问题在于：\textbf{你无法向 LLM 询问它是否有偏见}——它不具备可靠的自我审视能力。后来在对话中，聊天机器人确实修正了这个错误，但最初的输出已经暴露了潜在的偏见。
\end{warningbox}

\subsection{偏见不能通过禁词解决}

讲者指出一个常见误解：人们以为可以通过阻止某些词语来消除 LLM 的偏见。例如，阻止种族相关词汇出现在输出中。但偏见远比特定词汇更深层——它可以隐含在句子结构、话题选择、细节取舍等各个层面。

经典案例：自动皂液器由于传感器只针对浅色皮肤调校，导致深色皮肤用户无法触发出液。这不涉及任何``偏见词汇''，而是技术设计本身的系统性偏见。在大语言模型中，类似的偏见以更加隐蔽的形式存在于生成的代码和文本之中。

\subsection{本章小结}

AI 偏见比表面上看到的更加微妙和根深蒂固。它可能表现为代码中的父系优先、推荐结果中的人口统计偏差、或语言生成中的刻板印象。简单的禁词策略无法解决这个问题，需要系统性的测试和评估才能发现并缓解偏见。

